\section{威胁模型}
\label{sec:threat}

\paragraph{系统模型} 分诊智能体运行在组织内网。对每条告警，它收到一组候选原因（其中之一正在生效，可能是显式的 \textit{other}），以及一个覆盖安全遥测的只读探针目录：认证与 HTTP 日志、来源信誉、DNS 日志、配置、资产清单、变更工单和威胁情报。每个探针返回结构化记录或原始日志文本。智能体以一个结论（一个原因及其引用的观测）、一次弃权或预算耗尽结束。良性结论建议不采取行动而关闭案件；其他任何结局都把案件交给分析员。智能体从不执行处置，也从不向被调查系统发送流量。

\paragraph{攻击者能力与目标} 攻击者引发了告警，希望它被结为良性，至少不产生结论。攻击者无法接触智能体、其宿主或配置，只能通过智能体读取的证据施加影响：
\begin{description}\setlength\itemsep{0pt}
  \item[A1：指令注入。] 攻击者往会被复制进日志的字段（用户代理、请求载荷、命令行、工单评论）写入自由文本。只要模型看到原始探针响应，这段文本就会到达模型。
  \item[A2：证据伪造。] 攻击者控制某一个数据来源报告的内容，例如让流量经过看起来像扫描器的地址，或者污染几个探针共同读取的上游信誉源。
  \item[A3：读数操纵。] 当模型把原始日志解析为结构化结果时，攻击者往智能体可能读取的每条日志的自由文本字段写入内容，声称良性读数或伪造整条日志记录。攻击者知道探针目录、结果词表和防御方式。
\end{description}

\paragraph{信任假设} 我们信任控制器、预测表、探针实现和日志，攻击者无法修改它们。我们信任按文档日志格式编写的固定规则解析器，但不信任作为解析器的模型，也不信任模型会遵循合理的流程。环境会报告事件何时变化，预测表是在与待分诊案例相似的案例上估计的。

\paragraph{安全目标} \textbf{G1}：攻击者书写的文本（A1）或读数（A3）不能让任何攻击被结为良性。\textbf{G2}：伪造单个来源组（A2）不能让任何攻击被结为良性。\textbf{G3}：每个被接受的结论都引用了一条支持它的当前观测。\textbf{G4}：防御解决的诚实案例不少于最好的模型驱动智能体。

\paragraph{范围} 控制两个独立来源组或控制规则解析器的攻击者不在范围内；只拖延结论的拒绝服务攻击也不在范围内，因为未解决的案例会交给分析员。我们不生成新的候选原因，也不处理并发原因。
